\documentclass{article}
\usepackage[T1]{fontenc}
\usepackage{lmodern}
\usepackage{graphicx}
\usepackage{amsmath,amssymb}
\usepackage[a4paper,margin=2cm]{geometry}
\usepackage[absolute,overlay]{textpos}
\setlength{\TPHorizModule}{1mm}
\setlength{\TPVertModule}{1mm}
\usepackage[indonesian]{babel}
\usepackage{hyperref}
\setlength{\emergencystretch}{2em}
% unit: Halaman 20

\begin{document}

% === Halaman 20 (llm) ===

Metode enhanced-LSB bagus untuk citra dengan kontras tinggi, yaitu citra yang memiliki warna latar yang jelas atau memiliki perbedaan warna yang kontras antara latar dengan gambar utama

\vspace{20mm}

Gambar III-1 Gambar yang mengandung pesan rahasia dan hasil enhanced LSB-nya [PAU07]

% deskripsi: Perbandingan antara citra asli yang mengandung pesan rahasia dan hasil ekstraksi menggunakan metode enhanced LSB.
% bbox normalisasi: [0.2200, 0.5300, 0.7700, 0.7600]
\begin{textblock*}{115.50mm}(46.20mm,157.41mm)
\noindent\includegraphics[width=115.50mm,height=68.31mm]{../../images/llm/page_020_img_01.png}%
\end{textblock*}

\end{document}
